\section{Symmetry}

\begin{exercise}[Three vibrations of water]
\textbf{Experiment.}
Infrared spectroscopy finds three fundamental vibrational modes of a water
molecule: one bend and two stretches \cite{shimanouchi}.  The equilibrium
shape is preserved by a reflection and by interchange of the two hydrogen
atoms.  Let the resulting finite group \(G\) act on the small-displacement
space \(V\), and suppose the dynamical operator \(D:V\to V\) commutes with
that action.

Without assigning coordinates to the atoms, show that every eigenspace of
\(D\) is a \(G\)-representation.  Decompose \(V\) into irreducible
representations by averaging an arbitrary inner product over \(G\).  Explain
why this decomposition separates the symmetric stretch, antisymmetric
stretch, and bend.  Use the transformation of the dipole moment to state which
components can absorb infrared radiation and thereby explain the observed
selection rule.
\end{exercise}

\begin{exercise}[Missing spectral lines]
\textbf{Experiment.}
Atomic spectra contain sharp lines, but many energy differences produce no
line.  Suppose a symmetry group \(G\) acts unitarily on the state space \(H\),
the Hamiltonian commutes with \(G\), and the radiation-coupling operator
transforms as a representation \(W\).

For initial and final irreducible subspaces \(V_i,V_f\subset H\), interpret a
transition amplitude as a \(G\)-equivariant map
\[
        V_i\longrightarrow V_f\otimes W.
\]
Use Schur's lemma to show that the amplitude vanishes unless \(V_i\) occurs in
\(V_f\otimes W\).  Apply this statement to rotational symmetry and an electric
dipole, obtaining the usual change in angular momentum.  Explain why a
vanishing intertwiner predicts an absent line rather than merely a weak one.
\end{exercise}

\begin{exercise}[Hydrogen degeneracy]
\textbf{Experiment.}
Ignoring fine structure and external fields, hydrogen energy levels with
different orbital angular momentum share the same energy.  Spectroscopy
organizes the bound-state energy by one principal integer \cite{hydrogen}.
Let \(L\) be angular momentum and \(A\) the conserved Runge--Lenz vector.

Starting from their coordinate-free commutator relations on a fixed negative
energy eigenspace, rescale \(A\) and form two commuting triples
\[
        J_\pm=\tfrac12(L\pm A').
\]
Show that each triple generates a copy of \(\mathfrak{su}(2)\).  Determine the
dimension of the resulting representation and recover the degeneracy \(n^2\).
Explain which extra conserved quantity accounts for degeneracy not forced by
ordinary spatial rotations, and which experimental perturbations destroy it.
\end{exercise}

\begin{exercise}[A full turn changes a sign]
\textbf{Experiment.}
Neutron interferometry detects a relative phase of \(\pi\) when one beam's
spin is rotated through \(2\pi\), and return after \(4\pi\)
\cite{hagiwara}.  Spatial directions are transformed by
\(\mathrm{SO}(3)\), but spin states form rays in a two-dimensional Hilbert
space.

Show that a continuous action on rays lifts to an ordinary unitary
representation of the double cover \(\mathrm{SU}(2)\).  Prove that the
nontrivial element above the identity rotation acts as \(-\id\) in the
spin-\(\tfrac12\) representation.  Explain why the sign is invisible on an
isolated ray yet becomes observable when the rotated beam interferes with an
unrotated beam.  Identify the obstruction to replacing this projective action
by an ordinary representation of \(\mathrm{SO}(3)\).
\end{exercise}

\begin{exercise}[Protected end states]
\textbf{Experiment.}
Open integer-spin antiferromagnetic chains exhibit low-energy spin-\(\tfrac12\)
degrees of freedom at their ends while the bulk remains gapped.  Model the
ground state by a tensor network whose internal space carries a projective
representation \(u_g\) of the preserved symmetry group:
\[
        u_g u_h=\omega(g,h)u_{gh}.
\]

Use associativity to derive the cocycle identity for \(\omega\).  Show that a
change \(u_g\mapsto\beta(g)u_g\) changes \(\omega\) by a coboundary but cannot
change its cohomology class.  Explain why a nontrivial class permits
half-integer edge behavior even though every bulk site has integer spin.
Argue why a continuous, symmetry-preserving, gap-preserving deformation
cannot remove the edge representation.
\end{exercise}

\begin{exercise}[Zeeman splitting]
\textbf{Experiment.}
A spectral line splits into several components when an atom is placed in a
magnetic field, and the pattern depends on the field direction.
Let an irreducible rotation representation describe an unperturbed energy
eigenspace, and let the field select the subgroup of rotations about its axis.

Restrict the representation to this subgroup and decompose it into
characters.  Explain why the perturbation can assign different energies to
the resulting weight spaces.  Recover the number and spacing of components in
the simplest orbital model.  Explain why removing the field restores the
degeneracy without requiring the eigenvectors to return to a preferred basis.
\end{exercise}

\begin{exercise}[Ammonia inversion]
\textbf{Experiment.}
The ammonia microwave spectrum contains a close doublet associated with the
nitrogen atom tunneling between two classically equivalent configurations.
Let the reflection exchanging the configurations act unitarily on the
low-energy state space and commute with the Hamiltonian.

Decompose the space into the two characters of the reflection group.  Explain
why the energy eigenstates are symmetric and antisymmetric combinations rather
than states localized in one configuration.  Relate the doublet separation to
the tunneling amplitude, and explain why a measurement localizing the nitrogen
prepares a superposition which subsequently oscillates.
\end{exercise}

\begin{exercise}[Parity violation]
\textbf{Experiment.}
In polarized cobalt-60 decay, electrons are emitted preferentially opposite
the nuclear spin.  Reflection reverses momentum but not angular momentum.

Construct the scalar made from the spin and emitted momentum and determine
its behavior under spatial reflection.  Explain why a reflection-invariant
decay law would force the angular distribution to have equal rates in the two
directions.  Show how the observed asymmetry rules out parity as a symmetry of
the weak interaction while leaving rotational covariance intact.  Distinguish
failure of a physical symmetry from a change of coordinates.
\end{exercise}

\begin{exercise}[Identical-particle multiplets]
\textbf{Experiment.}
Atomic term spectra contain multiplets with systematic absences determined by
whether the electrons occupy symmetric or antisymmetric joint states.
Let the permutation group act on the tensor power of the one-electron state
space.

Decompose the tensor power into symmetry types.  Explain why the fermionic
state must lie in the alternating subspace and how this constrains the allowed
combination of spatial and spin factors.  Derive the singlet--triplet
alternative for two electrons.  Explain the observed missing terms as absent
representation components, not as transitions too weak to detect.
\end{exercise}

\begin{exercise}[Raman polarization rules]
\textbf{Experiment.}
Changing the relative polarizations of incident and scattered light makes
some molecular Raman lines appear and others disappear.  Let a vibrational
mode transform in a representation \(V\), and let the polarizability transform
in the symmetric square of the spatial vector representation.

Express the scattering amplitude as an invariant trilinear form.  Determine
when such an invariant exists by decomposing the tensor product.  Explain the
observed polarization-dependent absences as symmetry constraints on the
amplitude, and show why rotating the laboratory axes cannot change whether
the invariant subspace exists.
\end{exercise}
